\documentclass[12pt]{article}

\usepackage{amsmath,amssymb,booktabs,array,hyperref}
\topmargin -.5cm
\textheight 21cm
\oddsidemargin -.125cm
\textwidth 17cm

\newcommand{\D}{\Delta}
\newcommand{\cO}{\mathcal{O}}
\newcommand{\la}{\langle}
\newcommand{\ra}{\rangle}
\newcommand{\cC}{\mathcal{C}}

\begin{document}
\sloppy

\begin{center}
{\Large\bf How \texttt{2504.12290} Could Quantize the Axion Frequency}
\end{center}

\section{Purpose}

This note investigates how the mechanism of
\href{https://arxiv.org/abs/2504.12290}{arXiv:2504.12290},
which recovers quantized fluctuation data in a compact flux background, could be adapted to the
massless RR axion mode on $AdS_5\times S^5$.

The point is not to claim a complete derivation of the axion frequency from the current superstring SFT
setup.  The point is narrower: isolate the exact place where the local large-$R$ analysis fails to fix
$\omega$, and explain why the extended-BV/global-isometry method of \texttt{2504.12290} is the natural
mechanism that could restore the missing quantization condition.

\section{The Current Gap in the Projective-Patch Analysis}

For the axion zero mode we have been using the projective-coordinate ansatz
\begin{equation}
\phi(X)=e^{-i\omega\tau(X)}\,f(s),
\qquad
s=\frac{r^2}{R^2},
\qquad
r^2=\sum_{m=1}^4 (X^m)^2.
\end{equation}
The exact $X$-coordinate scalar equation gives the radial ODE
\begin{equation}
4s(1+s)f''(s)+4(2+3s)f'(s)+\frac{\omega^2}{1+s}f(s)=0.
\label{radialode}
\end{equation}
Expanding around the center of AdS,
\begin{equation}
f(s)=\sum_{n=0}^{\infty} a_n s^n,
\end{equation}
one finds a perfectly regular local solution for every chosen $\omega$.  The coefficients are fixed
recursively by \eqref{radialode}, but $\omega$ itself is not quantized by this inner expansion.

This is not an accident.  The expansion at fixed $X^a$ probes only the region
\begin{equation}
s=\frac{r^2}{R^2}\ll 1,
\end{equation}
i.e.\ a local patch near the center of AdS.  The actual quantization condition sits at the other end of
the radial problem, namely the AdS boundary $s\to\infty$, where one must require the absence of the
non-normalizable branch.

So the missing ingredient is global information, not another local recursion relation.

\section{What \texttt{2504.12290} Adds}

The essential lesson of \texttt{2504.12290} is that one should not study only one perturbative patch.
Instead, one introduces a family of Riemann normal coordinates centered at different target-space points,
together with a corresponding family of string fields related by diffeomorphism gauge transformations.
The background equation and these transport equations are then unified into an extended BV system.

The conceptual payoff is exactly the one we need here:
\begin{itemize}
\item the local weak-curvature expansion still controls each individual patch;
\item but the extended system transports the solution over the whole target space;
\item therefore the fluctuation problem can be organized globally under the target-space isometries.
\end{itemize}

Section 4.4 of \texttt{2504.12290} is the key point.  There the fluctuation string fields are required
to satisfy both:
\begin{itemize}
\item an on-shell equation,
\item and a covariance condition under the global isometry generators.
\end{itemize}
The fluctuations are organized into definite isometry representations.  Applying the covariance condition
twice identifies the target-space Laplacian with the quadratic Casimir of the isometry group, to the
relevant order in the weak-curvature expansion.  Matching to the target-space Laplacian spectrum then
recovers the allowed fluctuation representations.

The paper emphasizes that this works only because the extended fields know about the \emph{entire target
space}, not just the patch seen by expanding in $x/R$.

\section{\texorpdfstring{The AdS$_5$ Axion Analogue}{The AdS5 Axion Analogue}}

For the axion zero mode on $S^5$, the sphere representation is trivial.  So the only nontrivial
quantization problem is on the AdS side.

The analogue of the compact-space angular momentum label is now the lowest-weight representation of
$SO(2,4)$.  For a scalar primary one has the module
\begin{equation}
D(\D;0,0),
\end{equation}
whose quadratic Casimir is
\begin{equation}
\cC_2\big(D(\D;0,0)\big)=\D(\D-4).
\label{casimir}
\end{equation}
For the massless axion zero mode, the five-dimensional mass is
\begin{equation}
m_5^2=0.
\end{equation}
The standard AdS scalar equation is therefore
\begin{equation}
\left(\Box_{AdS}-m_5^2\right)\phi=0
\qquad \Longrightarrow \qquad
\D(\D-4)=0.
\label{massdim}
\end{equation}
So the two algebraic possibilities are
\begin{equation}
\D=0,\qquad \D=4.
\end{equation}

At the level of the local projective patch, both branches are visible.  This is exactly what the local
series solution of \eqref{radialode} is telling us: regularity near the center does not yet choose one
module over the other.

The global AdS state condition is what removes the ambiguity.  A physical bulk state must sit in the
normalizable lowest-weight module, not in the source branch.  For a massless scalar this means
\begin{equation}
\text{choose } D(4;0,0)\ \text{rather than}\ D(0;0,0).
\end{equation}
For the primary, the global time dependence is
\begin{equation}
e^{-i\omega\tau},
\qquad
\omega=\D,
\end{equation}
so the physical choice is
\begin{equation}
\omega=4.
\end{equation}

In this form the logic is parallel to the compact example of \texttt{2504.12290}:
\begin{itemize}
\item local perturbation theory gives differential equations but not the full discrete spectrum;
\item global covariance under the isometry group upgrades the fluctuation problem to a Casimir problem;
\item the physical spectrum is then obtained by selecting the appropriate global representation.
\end{itemize}

\section{Why the AdS Case Is Not Identical to the Sphere Case}

There is one important difference.  On a compact space the isometry representations are finite-dimensional,
and the quantization problem is directly the quantization of angular momentum or KK level.  In AdS the
relevant group is noncompact, so the representation data are lowest-weight modules rather than
finite-dimensional multiplets.

So the AdS analogue of the compact KK quantization condition is:
\begin{equation}
\text{normalizable lowest-weight } SO(2,4)\text{ module}.
\end{equation}
This is not enforced by the local $x/R$ expansion.  It is a global statement, exactly of the type that
the extended-BV formalism is designed to capture.

\begin{center}
\begin{tabular}{>{\raggedright\arraybackslash}p{0.24\textwidth}>{\raggedright\arraybackslash}p{0.31\textwidth}>{\raggedright\arraybackslash}p{0.31\textwidth}}
\toprule
 & \texttt{2504.12290} compact example & AdS$_5$ axion analogue \\
\midrule
Local patch data &
weak-curvature expansion in one Riemann normal patch &
large-$R$ projective expansion near the AdS center \\
Global symmetry input &
compact isometry representation / KK angular momentum &
$SO(2,4)$ lowest-weight module $D(\D;0,0)$ \\
Casimir relation &
target-space Laplacian matches compact-space Casimir &
$\Box_{AdS}$ matches $\cC_2=\D(\D-4)$ \\
What gets quantized &
allowed KK representation &
allowed normalizable global-energy branch \\
Physical selection rule &
representation appears in BRST cohomology globally &
choose the normalizable module, hence $\D=\omega=4$ \\
\bottomrule
\end{tabular}
\end{center}

\section{What This Suggests for the Current SFT Program}

The immediate lesson is not a new local calculation.  It is a change in what one should expect from the
existing local calculation.

The projective large-$R$ analysis around one origin should \emph{not} be expected to quantize $\omega$ by
itself.  Its role is to provide the inner solution.  The missing quantization condition must come from the
global completion of the fluctuation problem.

For the axion, that global completion is simple enough that one already knows the answer:
\begin{equation}
m_5^2=0,\qquad \cC_2=\D(\D-4),\qquad \text{normalizable branch } \Rightarrow \omega=\D=4.
\end{equation}
So the main value of the \texttt{2504.12290} perspective is not to discover a new number, but to explain
where that number should come from in SFT: from a globally covariant fluctuation problem, not from the
inner weak-curvature patch alone.

For Konishi the same mechanism should be harder but structurally similar.  One would need:
\begin{itemize}
\item the superstring analogue of the extended BV transport system on $AdS_5\times S^5$,
\item a globally covariant fluctuation equation for the relevant first-massive fields,
\item and a match between that equation and the appropriate $SO(2,4)\times SO(6)$ Casimirs.
\end{itemize}
Only at that stage should one expect a genuine quantization condition on the $O(1)$ part of the energy.

\section{Bottom Line}

The technique of \texttt{2504.12290} suggests the following resolution of the present puzzle.

\begin{equation}
\text{local large-}R\text{ patch} \neq \text{full frequency quantization problem}.
\end{equation}
The local patch determines the inner expansion of the axion mode, but the frequency is selected only when
the fluctuation problem is promoted to a global isometry/Casimir problem.  For the massless scalar axion,
that global statement is precisely the choice of the normalizable lowest-weight $SO(2,4)$ module
$D(4;0,0)$, which yields
\begin{equation}
\omega=4.
\end{equation}

\begin{thebibliography}{99}

\bibitem{Diff2504}
M.~Cho, S.~Collier, J.~Sandor and X.~Yin,
``Diffeomorphism in Closed String Field Theory,''
\href{https://arxiv.org/abs/2504.12290}{arXiv:2504.12290}.

\bibitem{SFT2507}
R.~Pius, B.~T.~Stefa\'nski and A.~A.~Tseytlin,
``Type IIB string field theory on $AdS_5\times S^5$ at large radius,''
\href{https://arxiv.org/abs/2507.12921}{arXiv:2507.12921}.

\bibitem{Witten}
E.~Witten,
``Anti-de Sitter space and holography,''
Adv.\ Theor.\ Math.\ Phys.\ {\bf 2} (1998) 253,
\href{https://arxiv.org/abs/hep-th/9802150}{hep-th/9802150}.

\bibitem{AharonyReview}
O.~Aharony, S.~S.~Gubser, J.~M.~Maldacena, H.~Ooguri and Y.~Oz,
``Large $N$ field theories, string theory and gravity,''
Phys.\ Rept.\ {\bf 323} (2000) 183,
\href{https://arxiv.org/abs/hep-th/9905111}{hep-th/9905111}.

\end{thebibliography}

\end{document}
